\definecolor{headerbg}{HTML}{d5f5e3}
\setlength{\arrayrulewidth}{1pt}

Les story points (SP) reflètent la complexité et l'effort de chaque tâche, estimés selon la suite de Fibonacci.

\begin{longtable}{|>{\raggedright\arraybackslash}p{4.2cm}|>{\raggedright\arraybackslash}p{4.8cm}|>{\raggedright\arraybackslash}p{4.5cm}|c|}
\hline
\rowcolor{headerbg}
\textbf{User Story} & \textbf{Tâche} & \textbf{Critères d'acceptation} & \textbf{\rotatebox{90}{SP}} \\
\hline
\endfirsthead
\hline
\rowcolor{headerbg}
\textbf{User Story} & \textbf{Tâche} & \textbf{Critères d'acceptation} & \textbf{\rotatebox{90}{SP}} \\
\hline
\endhead
\hline
\endfoot

%------ US-039 ------
\multirow{2}{=}{En tant que Chauffeur, je veux démarrer une livraison depuis l'application mobile et voir ses détails afin de préparer ma tournée.}
& Permettre au Chauffeur de consulter les détails d'une livraison (adresse, articles, créneau horaire) avant de la démarrer. & Les détails de la livraison s'affichent sur l'application mobile (adresse, liste d'articles, créneau horaire). & 2 \\
\cline{2-4}
& Mettre à jour automatiquement le statut de la commande au démarrage de la livraison et notifier le Planificateur. & Le statut passe à « En cours de livraison » ; le Planificateur est notifié en temps réel. & 1 \\
\hline

%------ US-040 ------
\multirow{2}{=}{En tant que Chauffeur, je veux confirmer une livraison avec la signature électronique du client depuis l'application mobile.}
& Permettre au Chauffeur de recueillir la signature du client sur écran tactile lors de la remise de commande. & La signature est capturée sur écran tactile, horodatée et transmise au système. & 3 \\
\cline{2-4}
& Faire passer la livraison au statut « Livrée » lors de la confirmation avec signature et notifier les parties concernées. & La livraison passe au statut « Livrée » ; la signature est rattachée à la livraison et consultable. & 2 \\
\hline

%------ US-041 ------
\multirow{2}{=}{En tant que Chauffeur, je veux prendre des photos géolocalisées comme preuve de livraison depuis l'application mobile.}
& Permettre au Chauffeur de prendre des photos géolocalisées lors de chaque livraison. & Les photos sont géolocalisées avec horodatage ; plusieurs photos peuvent être prises par livraison. & 2 \\
\cline{2-4}
& Rattacher les photos prises à la livraison concernée et les rendre consultables par le REC et le Manager. & Les photos sont consultables par le REC et le Manager dans l'historique de la livraison. & 1 \\
\hline

%------ US-042 ------
\multirow{2}{=}{En tant que Chauffeur, je veux valider les articles livrés depuis l'application mobile afin de confirmer la conformité de la commande.}
& Permettre au Chauffeur de valider les articles remis au client et détecter tout écart avec la commande prévue. & La validation des articles est réalisée ; tout écart entre les articles prévus et livrés génère une alerte. & 2 \\
\cline{2-4}
& Notifier le REC de tout écart constaté entre les articles prévus et effectivement livrés. & Un écart entre prévu et livré génère une alerte visible par le REC. & 1 \\
\hline

%------ US-043 ------
\multirow{2}{=}{En tant que Chauffeur, je veux signaler un problème de livraison depuis l'application mobile.}
& Permettre au Chauffeur de sélectionner le type de problème rencontré parmi une liste de motifs prédéfinis. & Le type de problème est sélectionné ; commentaire optionnel ; le statut de la livraison passe à « Échouée ». & 2 \\
\cline{2-4}
& Notifier automatiquement le Planificateur et le REC après le signalement d'un problème par le Chauffeur. & Le Planificateur et le REC reçoivent une notification immédiate avec les détails du problème. & 1 \\
\hline

%------ US-044 ------
\multirow{2}{=}{En tant que Planificateur, je veux suivre en temps réel la position des chauffeurs dans mon tableau de bord afin de superviser les tournées.}
& Afficher au Planificateur la position en temps réel de tous les chauffeurs actifs sur une carte. & La carte affiche tous les chauffeurs actifs ; les positions sont mises à jour régulièrement. & 3 \\
\cline{2-4}
& Afficher les alertes de retard directement sur la carte de suivi avec possibilité de filtrer par zone. & Les alertes de retard s'affichent automatiquement ; le filtrage par zone est disponible. & 2 \\
\hline

%------ US-045 ------
\multirow{2}{=}{En tant que Planificateur, je veux que le statut de la commande se mette à jour automatiquement à chaque action du Chauffeur afin de suivre l'avancement en temps réel.}
& Déclencher automatiquement les transitions de statut de la commande à partir des actions du Chauffeur sur l'application mobile. & Chaque action du Chauffeur déclenche la mise à jour du statut ; synchronisation en temps réel. & 2 \\
\cline{2-4}
& Notifier les parties prenantes concernées (REC, Client) à chaque changement de statut de la commande. & Les notifications sont envoyées aux bonnes parties prenantes à chaque transition de statut. & 1 \\
\hline

%------ US-046 ------
\multirow{2}{=}{En tant que REC, je veux générer un bordereau de livraison pour chaque tournée depuis mon application web.}
& Générer un document récapitulatif de tournée (bordereau) listant les livraisons, les articles et les clients concernés. & Le bordereau contient la liste des livraisons, les articles et les clients ; téléchargeable. & 2 \\
\cline{2-4}
& Permettre au REC de télécharger le bordereau depuis l'interface de gestion des tournées en un clic. & Le bordereau est généré et téléchargeable depuis l'interface en un clic. & 1 \\
\hline

%------ US-047 ------
\multirow{2}{=}{En tant que Chauffeur, je veux remplir un checklist de départ du véhicule sur l'application mobile avant le début de la journée.}
& Permettre au Chauffeur de remplir une liste de contrôles de départ avec prise de photos pour les anomalies constatées. & Le checklist est interactif ; les anomalies sont documentées avec photos. & 2 \\
\cline{2-4}
& Bloquer le démarrage de la tournée si les points critiques du checklist ne sont pas validés par le Chauffeur. & Le démarrage est bloqué si des points critiques ne sont pas validés ; une alerte est affichée. & 1 \\
\hline

%------ US-048 ------
\multirow{2}{=}{En tant que Client, je veux suivre ma livraison en temps réel sur une carte dans l'application mobile.}
& Afficher au Client la position du chauffeur sur une carte avec estimation de l'heure d'arrivée. & Le Client voit la position du chauffeur sur la carte ; l'estimation d'arrivée est mise à jour régulièrement. & 3 \\
\cline{2-4}
& Envoyer une notification au Client lorsque le chauffeur est à moins de 10 minutes de son adresse. & La notification est envoyée lorsque le chauffeur est à moins de 10 minutes de l'adresse. & 2 \\
\hline

%------ US-049 ------
\multirow{2}{=}{En tant que Planificateur, je veux que le système détecte automatiquement les retards et l'alerte afin de prendre des mesures correctives.}
& Détecter automatiquement un retard lorsqu'une livraison dépasse de plus de 15 minutes l'estimation initiale. & Une alerte est générée si le retard dépasse 15 minutes ; le statut de retard est affiché. & 2 \\
\cline{2-4}
& Notifier le Planificateur et le REC en cas de retard détecté avec les détails de la livraison concernée. & Le Planificateur et le REC reçoivent une alerte dans les 2 minutes suivant la détection. & 1 \\
\hline

%------ Mode 1 - Prépayé ------
\multirow{2}{=}{En tant que Planificateur, je veux voir le statut de paiement d'une commande afin d'autoriser ou bloquer la livraison.}
& Importer le statut de paiement depuis le système de gestion des commandes et l'afficher sur la fiche commande. & Le statut de paiement est importé avec la commande ; si prépayé, la livraison est autorisée automatiquement. & 2 \\
\cline{2-4}
& Afficher un indicateur « Prépayé » sur la commande visible par le Planificateur et le Chauffeur. & L'indicateur de paiement est visible sur l'interface web (Planificateur) et l'application mobile (Chauffeur). & 1 \\
\hline

%------ Mode 2 - Chèque ------
\multirow{2}{=}{En tant que Chauffeur, je veux photographier le chèque du client et l'envoyer au REC pour validation avant d'autoriser la livraison.}
& Permettre au Chauffeur de photographier le chèque et de l'envoyer en temps réel au REC pour traitement. & La photo est transmise et reçue par le REC en temps réel ; les informations du chèque sont extraites automatiquement. & 3 \\
\cline{2-4}
& Permettre au REC de valider le chèque et d'envoyer une autorisation de livraison au Chauffeur. & La validation par le REC déclenche une autorisation envoyée au Chauffeur ; en cas d'échec, le motif est communiqué. & 5 \\
\hline

%------ Mode 3 - Traite ------
\multirow{2}{=}{En tant que Chauffeur, je veux photographier une traite et l'envoyer au REC pour décision d'acceptation ou refus.}
& Permettre au Chauffeur de photographier la traite et de l'envoyer au REC pour décision. & La photo est transmise et reçue par le REC ; l'acceptation autorise la livraison ; le refus exige un motif. & 3 \\
\cline{2-4}
& Notifier le Chauffeur de la décision du REC (acceptation = livraison autorisée ; refus = retour marchandise). & Acceptation : livraison autorisée. Refus : motif enregistré et Chauffeur et Client notifiés. & 2 \\
\hline

%------ US-050 ------
\multirow{2}{=}{En tant que Manager, je veux consulter les indicateurs de livraison du jour dans mon tableau de bord.}
& Calculer et exposer au Manager les indicateurs de livraison du jour (livrées, en cours, en retard, échouées). & Le tableau de bord affiche les compteurs par statut mis à jour en temps réel. & 2 \\
\cline{2-4}
& Afficher au Manager les indicateurs de livraison du jour avec graphique d'évolution. & Les indicateurs s'actualisent automatiquement ; le taux de réussite et le graphique d'évolution sont affichés. & 1 \\
\hline

%------ US-063-067 ------
\multirow{2}{=}{En tant que Chauffeur, je veux consulter le récapitulatif de mes encaissements du jour depuis l'application mobile.}
& Calculer et exposer au Chauffeur le récapitulatif de ses encaissements du jour par mode de paiement et par montant. & Le récapitulatif affiche le total par mode (espèces, chèques, traites) et le nombre de transactions. & 2 \\
\cline{2-4}
& Afficher au Chauffeur l'écran d'encaissements avec liste détaillée et totaux par mode de paiement. & L'écran est accessible depuis l'application mobile ; les totaux sont corrects et mis à jour. & 1 \\
\hline

%------ US-068 ------
\multirow{2}{=}{En tant que REC, je veux rapprocher les encaissements avec les commandes livrées dans mon application web.}
& Permettre au REC de rapprocher les encaissements reçus avec les commandes livrées et de détecter les écarts. & L'association encaissement/commande fonctionne ; les écarts sont signalés et mis en évidence. & 3 \\
\cline{2-4}
& Afficher au REC la liste des encaissements avec indicateurs d'écarts et statut de rapprochement. & Le statut « Rapproché » est visible ; les écarts sont mis en évidence pour action corrective. & 2 \\
\hline

%------ US-105/107 ------
\multirow{2}{=}{En tant que Manager, je veux un tableau de bord stratégique avec export des rapports et visualisation des tendances.}
& Calculer les indicateurs stratégiques (taux de livraison, coût par kilomètre, satisfaction client, utilisation de la flotte) et les exposer au Manager. & Les indicateurs stratégiques sont disponibles avec comparaison par période. & 3 \\
\cline{2-4}
& Permettre au Manager d'exporter les rapports et de visualiser les tendances sur 30, 60 et 90 jours. & Les rapports sont exportables ; les graphiques de tendances sur 30/60/90 jours sont disponibles. & 2 \\
\hline


%------ US-105 ------
\multirow{2}{=}{En tant que Client, je veux soumettre une réclamation depuis l'application mobile afin d'obtenir un suivi.}
& Permettre au Client de soumettre une réclamation en précisant la catégorie, la description et des photos jointes. & La réclamation est créée avec catégorie, description et pièces jointes ; un numéro de suivi est généré. & 3 \\
\cline{2-4}
& Affecter automatiquement la réclamation au Responsable en charge du client concerné. & La réclamation est assignée au Responsable du client ; un numéro de ticket est communiqué au Client. & 2 \\
\hline

%------ US-106 ------
\multirow{2}{=}{En tant que Responsable des Engagements Clients, je veux voir la liste des réclamations et les traiter dans le respect du délai de service.}
& Mettre à disposition du REC la liste des réclamations avec filtres par statut, date et délai de service restant. & Les réclamations sont filtrables ; le délai de service restant (en jours) est affiché pour chaque réclamation. & 3 \\
\cline{2-4}
& Permettre au REC de prendre en charge, résoudre ou clôturer une réclamation avec notification au Client à chaque changement. & Le REC peut changer le statut ; le Client est notifié à chaque changement de statut. & 2 \\
\hline

%------ US-107 ------
\multirow{2}{=}{En tant que Client, je veux suivre le statut de ma réclamation depuis l'application mobile afin de rester informé de son traitement.}
& Permettre au Client de consulter l'état de sa réclamation avec l'historique des transitions de statut. & Le Client voit le statut en temps réel avec l'historique des transitions. & 2 \\
\cline{2-4}
& Notifier le Client à chaque changement de statut de sa réclamation. & Le Client est notifié en temps réel à chaque changement de statut. & 1 \\
\hline

%------ US-108/110 ------
\multirow{2}{=}{En tant que Manager, je veux des statistiques sur les réclamations et l'identification des causes récurrentes dans mon tableau de bord.}
& Calculer les indicateurs de réclamations (volume, délai moyen de traitement, taux de récurrence, causes principales) et les exposer au Manager. & Les statistiques incluent : nombre par catégorie, délai moyen, taux de récurrence. & 3 \\
\cline{2-4}
& Afficher les causes les plus fréquentes de réclamations dans une vue comparative dans le tableau de bord Manager. & Les principales causes de réclamations s'affichent par ordre de fréquence décroissante. & 2 \\
\hline

%------ US-109 ------
\multirow{2}{=}{En tant que Responsable des Engagements Clients, je veux que les réclamations non traitées soient escaladées automatiquement après 48 heures afin de respecter les délais de réponse.}
& Déclencher automatiquement l'escalade d'une réclamation non traitée après 48 heures avec notification au Manager. & L'escalade est déclenchée après 48 heures ; le Manager est notifié avec les détails de la réclamation. & 2 \\
\cline{2-4}
& Mettre en évidence les réclamations escaladées dans la liste du REC avec un indicateur d'urgence visible. & Les réclamations escaladées affichent un indicateur d'urgence dans la liste. & 1 \\
\hline

%------ US-119/120 ------
\multirow{2}{=}{En tant que Manager, je veux personnaliser mon tableau de bord avec des widgets configurables et exporter les rapports.}
& Permettre au Manager de réorganiser les widgets de son tableau de bord et de sauvegarder sa configuration. & Les widgets sont déplaçables ; la configuration est sauvegardée par profil utilisateur. & 3 \\
\cline{2-4}
& Permettre au Manager d'exporter les rapports du tableau de bord dans des formats standards. & Les rapports sont exportables depuis tous les tableaux de bord ; la mise en page est professionnelle. & 2 \\
\hline

%------ US-140/141/142 ------
\multirow{2}{=}{En tant que Chauffeur, je veux utiliser l'application mobile en mode hors ligne et que mes actions se synchronisent automatiquement au retour du réseau.}
& Permettre au Chauffeur d'utiliser les fonctionnalités essentielles de l'application sans connexion réseau, avec mise en file d'attente des actions. & Les données essentielles sont disponibles hors ligne ; les actions sont mises en file d'attente et exécutées au retour du réseau. & 5 \\
\cline{2-4}
& Synchroniser automatiquement les données au retour du réseau avec gestion des conflits éventuels et indicateur de statut de connexion. & La synchronisation démarre automatiquement ; les conflits critiques sont signalés ; un indicateur de connexion est visible. & 3 \\
\hline

%------ US-121/122/123 ------
\multirow{2}{=}{En tant que Manager, je veux voir les tendances sur 30, 60 et 90 jours et un tableau de bord de performance des chauffeurs avec rapports automatisés.}
& Calculer les indicateurs historiques sur 30, 60 et 90 jours avec détection des anomalies et les présenter au Manager. & Les graphiques affichent les tendances ; les anomalies sont détectées et mises en évidence. & 3 \\
\cline{2-4}
& Envoyer automatiquement des rapports récapitulatifs hebdomadaires aux destinataires configurés. & Les rapports hebdomadaires sont envoyés automatiquement aux destinataires configurés. & 2 \\
\hline

\caption{Sprint Backlog du Sprint 4 --- Exécution des Livraisons, Paiement \& Réclamations}
\end{longtable}
\setlength{\arrayrulewidth}{0.4pt}
